% ============================================================
%  X. Conclusion
% ============================================================
\section{Conclusion}

This paper presented DCPAM, a dual-camera-based
point-to-axis measurement model that replaces the
traditional steel wire with a laser reference line.
%
The complete mathematical pipeline---from pixel
extraction through coordinate transformation to
point-to-line distance computation---was derived
and validated through error propagation analysis
and Monte Carlo simulation.

% TODO: 补充结论——总结实验结果、最终达到的精度、
%       与传统方法的对比优势、未来工作方向


% ============================================================
%  Acknowledgments
% ============================================================
\section*{Acknowledgments}

% TODO: 补充致谢信息


% ============================================================
%  References
% ============================================================
\begin{thebibliography}{10}
\bibliographystyle{IEEEtran}

\bibitem{piotrowski2006shaft}
J.~Piotrowski,
\textit{Shaft Alignment Handbook},
3rd~ed.\hskip 1em
Boca Raton, FL, USA: CRC Press, 2006.

% TODO: 补充更多参考文献：
%   - 激光对中相关文献
%   - 双目视觉测量文献
%   - 相机标定文献（Zhang's method, COLMAP）
%   - CNN 回归相关文献
%   - 误差传播理论文献
%   - 工业对中标准（ISO, GB）

\end{thebibliography}
